% 2026-11-16 站主問：「設計國家的人可以當賭徒嗎」。
% Law 796--799 頁主對話自文字層讀過；Hume 引句自 Gutenberg #36120 核對，原檔存 readings/W10/sources/Hume_Essays_Gutenberg36120.txt。
% Federalist No. 51 經 Law 註 38（737 頁）轉引，未另取原書。
\chapter{憲法設計者的風險態度}
\label{ch:30_designer}

\begin{qaquestion}
Law\ix{persons}{law@Law, David S.} 第五部末尾分別給法院、暴君、憲法設計者「教訓」。設計國家制度的人，能否以賭徒的態度行事？
\begin{enumerate}
\tightlist
\item Law 給法院的建議，帶有何種風險態度？
\item Law 給憲法設計者的建議，帶有何種風險態度？
\item 制度設計的傳統，對設計者的風險態度有何要求？
\item Law 的設計建議本身押在什麼上面？
\end{enumerate}
\end{qaquestion}

依 Hume 轉述的設計格言與 Law 自己在 790 頁所作的權衡，設計者對政府一方不得下注，應假定統治者只顧私利而依最壞情形設計。Law 允許下注的是法院，796 頁稱受理爭議案件為審慎的賭注；設計者把監督與協調集中於一個唯一的憲法法院，而第一個服從均衡又無從保證\pg{799 頁}，這一安排本身仍是設計者無法避免的賭注，輸掉的是全體人民的自由。

\section{法院的賭注}
\label{sec:30_designer-1}

Law 在 796 頁用了 gamble 這個字。已有服從紀錄的法院，如美國聯邦最高法院與德國聯邦憲法法院，受理爭議案件「has little to fear and perhaps even something to gain」〔幾無可懼，甚或有所得〕；新興民主國家新設的憲法法院則面臨判決不被服從、服從聲譽胎死腹中（stillborn）的風險。案件若裁判成功，即生日後服從的預期，「If those gains seem more than commensurate with the risks involved, adjudication becomes a prudent gamble」〔所得若足以抵償所冒風險，裁判便成為審慎的賭注〕；策略性的法院不同於膽怯的法院，「nothing ventured, nothing gained」〔不入虎穴，焉得虎子〕。\footcite[796]{law2009}

下注者是法院，賭注是法院自己的服從聲譽。賭輸的後果，Law 在 781 頁已寫明：判決公然不被服從而無後果，法院的協調能力隨之瓦解，如同人人知道大家都不理會的紅綠燈。\footcite[781]{law2009}法院的賭注因此不只押上法院自己：法院一旦失去協調能力，人民藉以監督政府的信號也一併失去。

\section{設計者的建議}
\label{sec:30_designer-2}

\begin{enumerate}
\tightlist
\item
  797 頁：憲法設計者若要降低專制的可能、強化人民統治，「should simply do the opposite of what a dictator would do」〔照獨裁者的作法反過來做即可〕；司法審查最有效的形式，是由一個唯一而高度顯著的憲法法院，發出清楚、易懂、可自由取得的信號。\footcite[797]{law2009}
\item
  798 頁：讓政府官員任命監督他們的法官，「arguably akin to allowing the foxes to appoint those who will guard the henhouse」〔可說等於讓狐狸任命看守雞舍者〕；Law 主張人民選舉法官的作法值得更多考慮，可配合終身職或不得連任的固定任期與俸給保障。\footcite[798]{law2009}
\item
  799 頁：最難的是建立第一個服從均衡\ix{terms}{01@\zhlidxcat{賽局理論}!均衡}。Law 以騎腳踏車為喻，動起來之後容易維持，起步不易；他承認沒有明顯或萬無一失的方法使人普遍相信憲法法院有效，只能藉聯邦制或立法與行政分立，讓法院有機會在政治部門之間的爭議中與有力的一方站在一起，累積服從的紀錄。\footcite[799]{law2009}
\end{enumerate}

\section{制度設計傳統中的風險態度}
\label{sec:30_designer-3}

Hume\ix{persons}{hume@Hume, David} 1741 年〈Of the Independency of Parliament〉〔論國會之獨立〕開篇轉述政治著作家的格言：「in contriving any system of government, and fixing the several checks and controls of the constitution, every man ought to be supposed a knave, and to have no other end, in all his actions, than private interest」〔設計任何政府體制、安排憲法的各項制衡時，應假定人人皆為無賴，其一切行動除私利外別無目的〕；否則人民對自由與財產的保障，只剩統治者的善意，「that is, we shall have no security at all」〔亦即毫無保障〕。\footcite{hume1741}Hume 接著承認這條格言在政治上成立、在事實上不成立，理由是人在公共角色中比在私人生活中更不誠實。Law 在 737 頁註 38 引 Federalist No.~51 的「If men were angels, no government would be necessary」〔人若是天使，便無須政府〕，Hobbestown 的槍手正依此假定而設。\footcite[737 n.38]{law2009}

這一傳統要求設計者依最壞的情形設計：不押注於統治者的品格，以制度使其私利與公益相合。Law 第五部的賽局沿用了同一種態度。人民厭惡暴政的風險\pg{790 頁}，寧可承受法院誤判違憲，也不願承受誤判合憲，理由是暴政的代價難以承受、法院的過度積極較易承受\pg{789--790 頁}。\footcite[789--790]{law2009}在政府一方，設計者不賭。

\section{Law 的設計建議所押的賭注}
\label{sec:30_designer-4}

\begin{enumerate}
\tightlist
\item
  設計者把監督與協調集中於一個唯一而顯著的法院\pg{797 頁}，是把風險集中在一點。暴君的對策恰好是分散：多個管轄重疊的法院\pg{796--797 頁}。Law 要設計者反其道而行，所得是信號清楚，所失是這個唯一的法院一旦遭到俘虜或賭輸服從聲譽，協調裝置全面失效，沒有第二個法院可以替代。按 Hume 的格言，設計者應假定法官同樣是只顧私利的無賴；Law 對政府適用這條格言，對法院則以「no obvious means of oppressing the people」〔無明顯壓迫人民的手段〕\pg{790 頁}為由放寬。
\item
  第一個服從均衡無從保證\pg{799 頁}，則任何新設憲法法院的制度，在通過時都押注於人民與政治部門日後的預期。設計者能做的，是 799 頁所列的權力分立與聯邦制，增加法院站在有力一方的機會；賭局本身無法避免。
\item
  法院在 796 頁下注，用的是自己的服從聲譽。設計者下注，輸掉的是全體人民的自由與日後世代的制度，賭徒自負盈虧的前提在此不成立。Hume 的格言與 Law 790 頁的權衡都以此為據：在不可承受的損失面前，設計者應當依最壞情形設計。Law 的建議把賭注從設計者移到法院，法院可以作審慎的賭注，前提是設計者已經為法院賭輸的情形另備制度，例如 798 頁的選任方式與俸給、任期保障。Law 就此只提出選任方式的建議，法院賭輸之後人民如何另行協調，本文未論。\footcite[790, 796--799]{law2009}
\end{enumerate}
